\documentclass[utf8,twocolumn]{ctexart}
\usepackage{ctex}
\usepackage{amsmath}
\usepackage{circuitikz}
\usepackage{tikz}
\usepackage{graphicx}
\usepackage{hyperref}
\usepackage{geometry}
\usepackage{pdfpages}
\usepackage{booktabs}
\usepackage{subfigure}
\usepackage{float}
\usepackage{amsmath}
\usepackage{amssymb}
\usepackage{mathrsfs}
\usepackage{multirow}
\usepackage{inputenc}
\usepackage{fancyhdr}
\usepackage{physics}

\hypersetup{
    colorlinks=true,
    linkcolor=black,
    filecolor=black,
    urlcolor=black,
    citecolor=black,
}
\geometry{a4paper, scale=0.8}
\pagestyle{fancy}
\fancyhf{}
\lhead{全息光栅}
\rhead{无37~董皓彧~2023010659}
\cfoot{---~~\thepage~~---}
\setlength{\columnsep}{20pt}

\title{全息光栅实验报告}
\author{无37~董皓彧~2023010659}
\date{\zhtoday}

\begin{document}

\maketitle
\thispagestyle{fancy}

\section{实验过程及数据分析}

\subsection{扩束平行光}

搭建扩束准直光路，使用25倍显微物镜作为扩束镜，大口径凸透镜作为准直镜，两透镜共焦放置。使用平行平晶检验平行度，观察到反射方向的干涉条纹较宽，如图~\ref{fig:parallel}~所示，说明输出光束具有良好的平行性。

\begin{figure}[H]
    \centering
    \includegraphics[width=0.9\linewidth]{images/1-平行光晶检测干涉条纹.jpg}
    \caption{平行平晶检测干涉条纹}
    \label{fig:parallel}
\end{figure}

\subsection{全息光栅干涉条纹}

搭建马赫-曾德（M-Z）干涉仪合光光路，按照设计目标 $U = 28$ 条/mm 调节光路。在显微镜下观察到清晰的等间距平行干涉条纹，如图~\ref{fig:fringes}~所示。条纹间距均匀，对比度良好，可用于制作全息光栅。

\begin{figure}[H]
    \centering
    \includegraphics[width=0.9\linewidth]{images/2-MZ干涉器后干涉条纹.jpg}
    \caption{M-Z 干涉仪输出干涉条纹}
    \label{fig:fringes}
\end{figure}

受限于实验时间，未调节两束光夹角为 $\theta = 1^\circ$，因此实际干涉条纹间距可能不准。

\subsection{全息光栅衍射图样}

由于本次实验时间限制，未能完成全息干板的曝光显影及衍射图样的观测。根据理论分析，使用 $d \approx 36.3$ \textmu m 的全息光栅，He-Ne 激光（$\lambda = 632.8$ nm）的一级衍射角为
\begin{equation}
    \alpha_1 = \arcsin\!\left(\frac{\lambda}{d}\right) = \arcsin\!\left(\frac{632.8}{36300}\right) \approx 1.0^\circ
\end{equation}

\onecolumn
\section{预习报告}
\includepdf[pages=-]{images/预习报告-全息光栅.pdf}

\end{document}
